\documentclass[10pt,a4paper]{amsart}
\usepackage[margin=19mm]{geometry}
\usepackage{amsmath,amssymb,mathtools}
\usepackage[scr=rsfs,cal=euler]{mathalfa}
\usepackage{xfrac}
\usepackage[unicode,hidelinks,bookmarks=false]{hyperref}
\newcommand{\ie}{i.e.\ }
\newcommand{\cf}{cf.\ }
\newcommand{\resp}{respectively\ }
\newcommand{\Subsection}{\subsection}
\providecommand{\nobreakdash}{\nobreak\hbox{-}}
\setlength{\emergencystretch}{2em}
\multlinegap0pt
\usepackage{bm}
\usepackage{bbm}
\usepackage{dsfont}
\usepackage{xspace}
\usepackage{cancel}
\usepackage{mathtools}
\usepackage{cleveref}
\usepackage{amsthm}
\makeatletter
\@ifundefined{lemma}{\newtheorem{lemma}{Lemma}}{}
\makeatother
\makeatletter
\newcommand{\K}{\mathbb K}
\expandafter\ifx\csname conjecture\endcsname\relax
\newenvironment{conjecture}{\begin{conjecture0}}{\end{conjecture0}}
\fi
\expandafter\ifx\csname corollary\endcsname\relax
\newenvironment{corollary}{\bigskip \begin{corollary0}}{\end{corollary0}}
\fi
\expandafter\ifx\csname definition\endcsname\relax
\newenvironment{definition}{ \begin{defn0}}{\end{defn0}}
\fi
\expandafter\ifx\csname example\endcsname\relax
\newenvironment{example}{ \begin{example0}\rm}{\end{example0}}
\fi
\expandafter\ifx\csname lemma\endcsname\relax
\newenvironment{lemma}{\bigskip \begin{lemma0}}{\end{lemma0}}
\fi
\expandafter\ifx\csname proposition\endcsname\relax
\newenvironment{proposition}{\bigskip \begin{prop0}}{\end{prop0}}
\fi
\expandafter\ifx\csname remark\endcsname\relax
\newenvironment{remark}{ \begin{remark0}\rm}{\end{remark0}}
\fi
\expandafter\ifx\csname theorem\endcsname\relax
\newenvironment{theorem}{\bigskip \begin{thm0}}{\end{thm0}}
\fi
\makeatother
\title[A refinement on the local cactus rank algorithm]{A refinement on the local cactus rank algorithm}
\author{Alessandra Bernardi, Oriol Reig Fit\'e}
\date{Source: arXiv:2508.15062v1 (CC BY 4.0)}
\begin{document}
\maketitle

\noindent\emph{Source and licence.} A refinement on the local cactus rank algorithm. Version arXiv:2508.15062v1 (CC BY 4.0), distributed under the Creative Commons Attribution 4.0 International licence, as recorded in the arXiv licence metadata for this exact version and shown on the abstract page https://arxiv.org/abs/2508.15062v1. Full rights notice: licenses/arxiv-2508.15062-CC-BY-4.0.txt.\par\smallskip

\noindent\textit{Editorial scope.} Counted unit: the Lemma whose source label is \texttt{lemma: BasisFromHF} in the article (source0.tex lines 1284--1286, source label \texttt{lemma: BasisFromHF}), delivered together with the article's own complete original proof of it (source0.tex lines 1288--1304, one \texttt{proof} environment of 17 lines). No mathematics is written, repaired or re-derived: statement and proof are the article's own text line for line. Required context retained uncounted: affine:tranlsation (lines 231--233); base:not:change (lines 1242--1246). Every definition, display and label of those lines is retained unchanged. Omitted from this package and not redistributed: 1--230, 234--1241, 1247--1283, 1287--1287, 1305--2619. Nothing inside a retained excerpt is omitted or altered: every retained body is delivered exactly as the article's lines are written, so no omission receipt of any kind accompanies this package. Citations: the retained excerpts carry no citation command at all, so no bibliography entry of the article is transcribed and no reference is redistributed. Reference adaptation: none was needed and none was made; every reference inside the delivered unit resolves inside it, so no unresolved reference or citation is produced. Printed slips kept literal: no printed word, symbol, hypothesis or number of the counted statement or of its proof was altered. Layout and class: the article's own class is \texttt{amsart} with packages including xy, xcolor, graphicx, amsmath, amscd, amsthm, amssymb, dsfont, MnSymbol, algorithm, algpseudocode, tikz-cd, tikz, relsize; the wrapper is the lane's pinned \texttt{amsart} editorial wrapper at 10pt on A4, which restates the theorem-like environments the retained text needs and adds the article's own macro definitions that the retained text needs (\texttt{\textbackslash K}), with extra wrapper packages (bm, bbm, dsfont, xspace, cancel, mathtools, cleveref). The delivered numbering is the unit's own. Assets: no retained span contains a figure environment, a graphics-inclusion command, a PDF-page inclusion, a table or a TikZ picture; no image file is copied into this package, the package declares no asset dependency and no image file is redistributed with it. Licence: version arXiv:2508.15062v1 of this article is distributed under the Creative Commons Attribution 4.0 International licence, as recorded in the arXiv licence metadata for this exact version and shown on the abstract page https://arxiv.org/abs/2508.15062v1. A full rights notice is delivered at licenses/arxiv-2508.15062-CC-BY-4.0.txt. Characters: every retained excerpt is pure ASCII, so the delivered engine, XeLaTeX with its default Latin Modern text font, reports no missing character.
\begin{equation}\label{affine:tranlsation}
\varphi: R \to R, \qquad \varphi(x_i) = x_i + \zeta_i
\end{equation}

\begin{lemma}\label{base:not:change}
Let $A=R/I$ be a local Artinian ring supported at $\zeta$ with a complete staircase basis $B$. After the affine change of coordinates $\varphi$ of \eqref{affine:tranlsation}, 
%sending $\zeta$ to $0$
 the set $B$ remains a basis of $\varphi(A)$.
\end{lemma}

\begin{lemma}\label{lemma: BasisFromHF}
    Let $A= \K[x_1 , \ldots , x_n]/I$ be an Artinian local ring of codimension $n$. Then there exists a complete staircase of monomials $B$ whose reduction modulo $I$ is a $\K$-basis, and the number of elements in $B$ of degree $i$ is $\operatorname{HF}_A(i)$.
\end{lemma}

\begin{proof}
    By \Cref{base:not:change} and the invariance of the Hilbert function of a local Artinian algebra under affine change of coordinates, we may assume that $A$ is supported at the origin.  
    We will prove the existence of a basis $B$ of $A$ consisting of monomials forming a complete staircase, with the degree distribution prescribed by the Hilbert function. The proof is in three steps:

    \begin{itemize}
        \item Existence of a monomial basis respecting the degree distribution.  
        Let $\mathfrak{m}$ be the maximal ideal of $A$. For each $i \ge 0$, the quotient $\mathfrak{m}^i / \mathfrak{m}^{i+1}$ is spanned by the residue classes of monomials of total degree $i$.  
        Choosing a $\K$-basis of $\mathfrak{m}^i / \mathfrak{m}^{i+1}$ made of monomials for each $i$, and collecting them for all $i$, we obtain a monomial basis $B$ of $A$. The number of basis elements of degree $i$ is then $\dim_{\K} \mathfrak{m}^i / \mathfrak{m}^{i+1} = \operatorname{HF}_A(i)$.

        \item Staircase property.  
        Suppose $b \in B$ has degree $i$ and $x_j \in \{x_1,\ldots,x_n\}$. If $b x_j$ is nonzero in $\mathfrak{m}^{i+1}/\mathfrak{m}^{i+2}$, then $b \notin \mathfrak{m}^{i+1}$ (otherwise $b x_j \in \mathfrak{m}^{i+2}$, contradicting $b x_j \neq 0$).  
        This shows that whenever $b \in B$ and $b x_j \neq 0$ in $A$, the product $b x_j$ lies in the basis of degree $i+1$. Hence $B$ can be chosen to be closed under multiplication by variables, i.e. a complete staircase.

        \item Inclusion of $1$ and the variables.
        Since $\operatorname{HF}_A(0) = 1$ and $\operatorname{HF}_A(1) = n$, the constant $1$ and the classes of $x_1,\ldots,x_n$ must appear in $B$.
    \end{itemize}
\end{proof}

\end{document}
